% !TeX root = ../report.tex

\phantomsection
\section*{PHỤ LỤC}
\addcontentsline{toc}{section}{PHỤ LỤC}

\subsection*{Phụ lục A. Dữ liệu thô đo độ trễ vòng lặp E2E (\texttt{latency\_results.csv})}

Bảng \ref{tab:raw-latency-csv} trích nguyên văn toàn bộ các trường dữ liệu thô thu được từ kịch bản đo hiệu năng \path{03_CoreIOT_Config/scripts/measure_latency.py} và lưu tại \path{04_Docs_Media/experiments/latency_results.csv}.

\begin{table}[H]
  \centering
  \caption{Dữ liệu đo độ trễ vòng lặp E2E trích từ \texttt{latency\_results.csv}}
  \label{tab:raw-latency-csv}
  {\footnotesize
  \setlength{\tabcolsep}{3pt}
  \renewcommand{\arraystretch}{1.25}
  \begin{tabularx}{\textwidth}{|C{0.05\textwidth}|C{0.18\textwidth}|C{0.09\textwidth}|C{0.09\textwidth}|C{0.08\textwidth}|C{0.09\textwidth}|C{0.08\textwidth}|Y|}
    \hline
    \rowcolor{iotTableHead}
    \textbf{Lượt} & \textbf{Timestamp (UTC)} & \textbf{PUBACK (ms)} & \textbf{First RPC (ms)} & \textbf{Servo (ms)} & \textbf{Buzzer (ms)} & \textbf{LED (ms)} & \textbf{All 3 RPC / Trạng thái} \\
    \hline
    1 & 2026-11-22 04:41:37 & 1.58 & 88.01 & 88.01 & 88.50 & 88.35 & 88.50 ms (\path{PASS_ALL_3}) \\
    \hline
    2 & 2026-11-22 04:41:39 & 2.39 & 113.02 & 113.26 & 113.02 & 113.43 & 113.43 ms (\path{PASS_ALL_3}) \\
    \hline
    \multicolumn{2}{|r|}{\textbf{Trung bình (Avg)}} & \textbf{1.99} & \textbf{100.52} & \textbf{100.64} & \textbf{100.76} & \textbf{100.89} & \textbf{100.97 ms (2/2 Đạt)} \\
    \hline
  \end{tabularx}
  }
\end{table}

\subsection*{Phụ lục B. Đặc tả REST API cục bộ trên Local IoT Gateway (Cổng 5050)}

Máy chủ HTTP nhúng \path{EmbeddedWebMonitor} trong \path{02_Local_Gateway/ui/gateway_infrastructure.py} cung cấp hai điểm cuối phục vụ giao tiếp giữa trình duyệt SCADA và tiến trình Gateway:
\begin{itemize}[leftmargin=*, itemsep=2pt]
  \item \textbf{\path{GET /api/telemetry}:} Trả về đối tượng JSON (\texttt{HTTP 200 OK}) chứa ảnh chụp trạng thái thời gian thực gồm các khóa: \path{smoke}, \path{flame}, \path{temperature}, \path{humidity}, \path{tinyml_risk}, \path{servo}, \path{buzzer}, \path{led}, \path{alarm_override}, \path{is_anomaly}, \path{anomaly_score}, \path{alarm_status}, \path{alarm_notification}, \path{last_seen}, \path{cloud_connected}, \path{buffer_count} và \path{gateway_type}.
  \item \textbf{\path{POST /api/control}:} Nhận thân yêu cầu JSON (\path{Content-Type: application/json}) chứa khóa điều khiển cưỡng bức (ví dụ \texttt{\{"buzzer": 1\}}, \texttt{\{"led": 1\}}, \texttt{\{"servo": 90\}}), đưa vào \path{command_queue} để gửi xuống UART ESP32-S3 và phản hồi \texttt{\{"status": "success", "command": \{...\}\}}.
\end{itemize}

\subsection*{Phụ lục C. Khai báo phạm vi sử dụng công cụ AI}

Thực hiện đúng quy định về liêm chính học thuật của môn học CO3054, nhóm khai báo phạm vi sử dụng các công cụ trí tuệ nhân tạo hỗ trợ trong quá trình thực hiện Bài tập lớn:
\begin{itemize}[leftmargin=*, itemsep=2pt]
  \item \textbf{Hỗ trợ lập trình và kiểm thử:} Sử dụng trợ lý lập trình AI để gợi ý cấu trúc mô-đun Python (hàng đợi bất đồng bộ, SQLite WAL, kịch bản kiểm thử đơn vị \texttt{unittest}), rà soát lỗi cú pháp C++/Arduino và cấu trúc JSON của ThingsBoard 3.x. Toàn bộ mã nguồn đều được thành viên nhóm trực tiếp kiểm tra, chỉnh sửa, nạp chạy trên bo mạch thật YoloUNO ESP32-S3 và chịu trách nhiệm giải trình khi vấn đáp.
  \item \textbf{Hỗ trợ trình bày báo cáo LaTeX:} Sử dụng công cụ AI để chuẩn hóa cú pháp LaTeX, kiểm tra escape ký tự kỹ thuật, tổng hợp số liệu từ mã nguồn thật và tệp đo đạc \path{latency_results.csv}. Mọi thông số phần cứng, sơ đồ chân, kết quả đo đạc và hình ảnh minh chứng đều được trích xuất trung thực từ sản phẩm thực tế của nhóm.
\end{itemize}
